\documentclass[11pt]{article}
\usepackage[margin=1in]{geometry}
\usepackage{amsmath,amssymb,microtype}
\usepackage[colorlinks=true,urlcolor=blue]{hyperref}
\title{Literature status: quantitative convergence in the inconsistent proximal-splitting case}
\author{Automated mathematics run \texttt{fa\_banach\_001}, lane 2}
\date{11 August 2026}
\begin{document}
\maketitle

\noindent\textbf{Status:} literature already answered for open problem 1,
in the form of quantitative sufficient conditions.

\section*{Source question}
Florian Lauster and D. Russell Luke,
\emph{Convergence of Proximal Splitting Algorithms in
$\operatorname{CAT}(\kappa)$ Spaces and Beyond}, arXiv:2105.05691,
Section 6, physical PDF page 17.  Its first open direction is to determine
requirements for quantitative convergence of proximal splitting methods when
the individual proximal mappings do not have common fixed points---the
``inconsistent case.''

\section*{Supporting answer}
D. Russell Luke and Mahshid Mirhashemi,
\emph{Quantitative Convergence of Proximal Splitting Iterations in Uniformly
Convex Metric Spaces}, arXiv:2605.03484.  The abstract says explicitly that
the theory does not assume common minima or vanishing proximal parameters.
The introduction (physical pages 1--2) cites the source as reference [29],
describes [29, Theorem 25] as a consistent-case convergence-rate result, and
identifies quantitative convergence for inconsistent instances as the missing
case addressed by the new paper.

Theorem 3.3 (physical page 9) treats cyclic and relaxed compositions
\[
 T=\operatorname{prox}^{p}_{f_m,\lambda_m}\circ\cdots\circ
   \operatorname{prox}^{p}_{f_1,\lambda_1}
\]
and their geodesically relaxed analogues on $p$-uniformly convex spaces with
curvature bounded above.  Under its explicit regularity, fixed-point
existence/invariance, and stability-gauge hypotheses, parts (iv)--(v) give
quantitative convergence and, for a linear gauge, $R$-linear convergence.
No intersection $\bigcap_j\arg\min f_j$ is assumed.  Corollary 3.4, beginning
on physical page 10, specializes these conclusions to Hadamard spaces and
removes the general composition assumptions in the parameter range stated
there.

\section*{Provenance and scope}
This identification is author-aware, not merely inferred: Luke coauthored
both papers, and the supporting paper explicitly cites and contrasts the
source theorem before announcing the removal of common minima.  It therefore
belongs in \texttt{literature\_already\_answered}.

The answer is scoped.  It supplies sufficient requirements, but not an
intrinsic necessary-and-sufficient characterization.  Existence/invariance
and a metric-subregularity-type stability gauge remain assumptions; in
positive curvature, additional composition/relaxation hypotheses remain.
The supporting paper's conclusion leaves removal of some of these assumptions
open.  The source's second open direction---automatic metric subregularity
from nonlinear structural hypotheses---is not resolved and is reiterated as
future work.

\section*{Evidence}
The packet contains the original arXiv PDF as \texttt{source\_paper.pdf} and
the decisive later PDF as
\texttt{supporting\_paper\_2605.03484.pdf}.  The source question is on physical
page 17; the supporting identification is on pages 1--2, with the main theorem
and Hadamard corollary on pages 9--10.

\begin{thebibliography}{9}
\bibitem{LL2021}
F. Lauster and D. R. Luke,
\emph{Convergence of Proximal Splitting Algorithms in
$\operatorname{CAT}(\kappa)$ Spaces and Beyond}, arXiv:2105.05691 (2021).

\bibitem{LM2026}
D. R. Luke and M. Mirhashemi,
\emph{Quantitative Convergence of Proximal Splitting Iterations in Uniformly
Convex Metric Spaces}, arXiv:2605.03484 (2026).
\end{thebibliography}

\end{document}
